\chapter{Discussion}

\begin{itemize}

    \item Hydraulic balancing has a significant impact on system performance; as discussed in 
    [Enhancing district heating…],  even a single consumer's valve can noticeably affect the entire network.

    \item The Bayesian Optimization approach introduces serious assumptions that warrant further scrutiny, 
    including the choice of acquisition function, kernel, and hyperparameter settings.

    \item The deviation between the model and the data from Van der Heijden may be partly caused by 
    the mass flow rate not being controllable with sufficient precision.

    \item The PhD thesis provides a relevant discussion on valve actuation control that could inform 
    further refinements [phdThesis].

    \item On page 54 of the PhD thesis, the valve actuator is described in detail. This level of 
    modelling is not applied here, leaving room for increased precision. The relevance depends on 
    the required accuracy. Additionally, hysteresis is noted to apply primarily to large valves, 
    not necessarily to small valves used in thermal end units.

    \item A constant Reynolds number is assumed, which is physically incorrect. This also renders the 
    Haaland equation invalid in cases where the actual Reynolds number falls below its applicable range.

    \item Minor losses are not accounted for in the hydraulic model.

    \item Scaling the UA value based on the Reynolds number could improve the realism of the heat 
    exchanger model, bringing it closer to physically observed behaviour.

    \item The Node method limits the ability to accurately control pressure across a pipeline, 
    though the extent of this limitation depends on the specific configuration.

    \item Valve opening displacement introduces a delay, though this may be negligible given the 
    one-minute time step used in the simulation.

    \item Valve hysteresis and other non-ideal valve characteristics are not incorporated into the model.

    \item The heat capacity in the Node method is overweighted because the pipe temperature $T_\text{pipe}$ 
    is not assumed to decrease, which may cause the model to cool down more slowly than reality, 
    particularly when convection and conduction interact differently at low or zero flow.

    \item The Node method is not well-suited for zero mass flow conditions, as pipe temperatures do 
    not decrease when water is stationary. This introduces a discrepancy between convective and 
    conductive heat transfer, potentially causing the model to underestimate cooling rates.

    \item Overshoot in heat supply may result from the PI controller being tuned for a single floor 
    rather than being adjusted on a per-floor basis.

    \item The NTU method assumes that the fluid passes through the heat exchanger within a single 
    time step and exchanges all energy instantaneously, whereas the NTU formulation fundamentally 
    assumes steady-state conditions. So you immediately assume that all the water within the heat exchanger is exchanged while you could also implement a time delay step. To more mimic this behavior like Femke did in here thesis p70.

    \item The pump model could be extended substantially to capture more realistic behaviour.

    \item Low flow rates in the model compared to reality may stem from overflow dynamics. In practice, 
    the valve opens significantly further than what is achievable in the current modelling setup, 
    at least within the available calibration time. See the elaboration in \textit{Obtain Realistic Parameters}.

    \item Thermal end unit keep-warm settings are not accounted for in the consumer model.

    \item All consumer units are modelled with identical assumptions, without accounting for the 
    fact that upper-floor apartments receive less flow than lower-floor ones.

    \item Directly adopting the pipe water temperature at small time steps may introduce errors, 
    as in reality the pipe does not lose energy as rapidly as the model assumes.

    \item Extending the valve model to include non-linearities such as hysteresis and deadband raises 
    the question of compatibility with Bayesian Optimization, which is generally not well-suited 
    for highly non-linear and discontinuous input-output relationships https://instrunexus.com/understanding-control-valve-hysteresis-deadband-response-time/.

    \item The consumer-side flow rate naturally decreases with height in the building, which is not 
    reflected in the current model.

    \item Assuming a constant ambient temperature is a simplification that does not reflect realistic 
    operating conditions.

    \item A more realistic model for supply temperature production and return temperature handling, 
    such as an ATES model, could be implemented to improve overall system representation.

    \item The supply temperature is assumed to be directly controllable, without modelling the 
    heat supply system itself.

    \item All assumptions made in selecting the hyperparameters of the Bayesian Optimization deserve 
    explicit justification and sensitivity analysis.

    \item too good working overflow valve which causes a low mass flow.It might be not realistic so this could have an important impact as there might be some minimal flow or sometime

    \item pressure drop over heat exchanger is not modelled detailed enough (like in the thesis of Femke Janssen). 

    \item mogelijk überhaupt meer aandacht besteden aan de dimensieloze getallen zoals Pr / Nu door de kleine tijdsstap als vervolg. 

    \item verschil tussen toestel wachttijd en leidingwachttijd. Bij mij wordt er vanuit gegaan dat het water meteen bij de klant is. En ik meet vanaf het toestel. Dus dat is meteen 35 seconden, niet tot aan de tap. 

\end{itemize}